\section{Idealismo y comercialización: quienes abren camino}

Por último, el texto vuelve a los juicios de valor: una empresa que cotiza en bolsa debe generar ganancias, pero Zhang Peng (张鹏) sigue anteponiendo la AGI y la contribución a la industria. En este punto toda la entrevista pasa de la historia de la empresa a su ideal de largo plazo.

Al final de la entrevista, le preguntaron a Zhang Peng qué elegiría entre construir una empresa que logre la AGI y construir una empresa muy rentable. Eligió la AGI, aunque también subrayó que ambas metas no se oponen: si se logra la AGI, también se tendrá una gran empresa en términos comerciales. Además, no estaría satisfecho con una Zhipu (智谱) que solo generara ganancias sin producir tecnología ni contribuir a la industria.

\begin{figure}[H]
\centering
\includegraphics[width=0.96\textwidth]{idealism-commercialization.png}
\caption{Idealismo y comercialización: ideal tecnológico, cultura de ingeniería, valor para el cliente, gobierno corporativo e identidad de pionero.}
\end{figure}

\begin{knowledgebox}{Lectura de la figura: Zhipu busca un camino equilibrado}
La figura va del ideal tecnológico al valor para el cliente y de ahí al gobierno corporativo tras la salida a bolsa, para volver al final a la idea de «pionero». Zhang Peng no insiste en proclamar la AGI sin más ni en ser solo una máquina de generar ganancias, sino en convertir «hacer que las máquinas piensen como los seres humanos» en un sistema de ingeniería capaz de servir a los clientes y a la sociedad.
\end{knowledgebox}

\subsection{Hacer que las máquinas piensen como los seres humanos}

El idealismo no puede quedarse en la palabra «AGI». Esta sección descompone el slogan de Zhipu para ver cómo apunta a la vez a un objetivo tecnológico y a un valor social.

Zhang Peng explica cómo entiende Zhipu la AGI a partir del slogan de la empresa: hacer que las máquinas piensen como los seres humanos y, a su vez, que fortalezcan a la sociedad humana. La frase tiene dos partes: la primera es un objetivo tecnológico, que las máquinas tengan mayor capacidad de comprensión, razonamiento y acción; la segunda es un objetivo de valor, que la tecnología genere un valor real en lugar de quedarse en el laboratorio.

\begin{importantbox}{Cómo volver operativa la AGI}
Si las definiciones de AGI difieren entre sí, para juzgar si una empresa persigue la AGI no basta con escuchar sus consignas: hay que ver en qué tareas por etapas, rutas tecnológicas, trayectorias de producto y valor para el cliente descompone su objetivo de largo plazo.
\end{importantbox}

\subsection{Pioneros, no solo proyectos estrella}

Lo anterior trata de «qué hacer»; aquí se trata de «cómo se quiere ser recordado». La palabra «pionero» es más sobria que «empresa estrella» y encaja mejor con el relato que Zhipu hace de sí misma.

Al inicio del programa se mencionó que hay quien compara a Zhipu con el «cemento»: hace un trabajo sólido, pero no despierta mucho entusiasmo. Al final, Zhang Peng dijo que espera que la historia recuerde a Zhipu como pionera de la AGI, como quien abrió camino. Este relato no contradice la imagen del «cemento»: quien abre camino no necesariamente es quien más atención atrae, pero tiene que pavimentar el camino durante mucho tiempo, tropezar con los obstáculos, convertir la investigación en aplicaciones, servir a los clientes y asumir las exigencias del gobierno corporativo.

\subsection{Resumen de la sección}

Zhipu se posiciona en un camino equilibrado entre el idealismo y la comercialización. La salida a bolsa no es el punto final, sino lo que permite a una empresa que persigue la AGI seguir avanzando bajo exigencias de gobierno corporativo más estrictas.
